\section{BPS sectors, black-hole diagnostics, and protected geometry}
\label{sec:bps-black-holes}

The motivation for a geometric formulation of eigenstate thermalization is
particularly sharp in a supersymmetric sector.  Ordinary spectral diagnostics
are built from the relative positions of distinct energy levels.  In a BPS
sector at fixed charges, however, supersymmetry fixes the energy through the
central-charge bound.  The protected states then form a degenerate eigenspace,
and the spacing distribution contains no information beyond the imposed
degeneracy.  This is a different obstruction from the failure of a finite-size
random-matrix fit: the relevant spacings vanish identically before any
statistical analysis is performed.  The standard ETH constructions of
Refs.~\cite{deutsch1991,srednicki1994} and the spectral-chaos criterion of
Ref.~\cite{bohigas1984} therefore do not, by themselves, provide a diagnostic
for this problem.

The natural object is the bundle of protected states over the space of
couplings.  Let $H(\lambda)$ be a family of Hamiltonians, let
$P(\lambda)$ project onto a fixed-charge BPS eigenspace of rank $D$, and write
$Q=1-P$.  When the eigenspace remains isolated from the non-BPS spectrum, its
off-fiber response is
\begin{equation}
 X_\mu=Q\,\partial_\mu P\,P,
 \qquad
 (X_\mu)_{ma}=\frac{\langle m|\partial_\mu H|a\rangle}{E_{\mathrm{BPS}}-E_m},
 \label{eq:bps-projector-response}
\end{equation}
where $|a\rangle$ belongs to the protected fiber and $|m\rangle$ belongs to
its complement.  The connection on the fiber and its curvature are
\begin{equation}
 (A_\mu)_{ab}=i\langle a|\partial_\mu b\rangle,
 \qquad
 F_{\mu\nu}=\partial_\mu A_\nu-\partial_\nu A_\mu-i[A_\mu,A_\nu].
 \label{eq:bps-berry-curvature}
\end{equation}
Under a change of frame in the degenerate fiber, $F_{\mu\nu}$ transforms by
conjugation.  Its eigenvalues, traces of powers, and holonomies are therefore
the appropriate gauge-covariant observables.  The construction is the
degenerate-state extension of the Berry phase and the Wilczek--Zee connection
\cite{berry1984,wilczekzee1984}; it is also the object proposed specifically
for BPS chaos in Ref.~\cite{chen2024bps}.

\subsection{What is known about smooth and black-hole-motivated sectors}

The distinction between smooth horizonless states and black-hole microstates
is a statement about the geometry represented by a particular family of
states, not a consequence of BPS protection alone.  D-brane counting and the
AdS/CFT correspondence provide the setting in which supersymmetric black-hole
microstates can be compared with states represented by smooth geometries
\cite{stromingervafa1996,maldacena1998ads}.  They do not imply that every
protected multiplet has the same response statistics.

The direct Berry-curvature comparison is supplied by
Ref.~\cite{chen2026berry}.  For the studied $1/2$-BPS states in the D1/D5 CFT,
which are used there as smooth, horizonless examples, the curvature under
marginal deformations is non-random and is often exactly zero at generic
coupling.  The qualification ``often'' is essential: the result concerns the
sectors and deformations analyzed in that work, not every D1/D5 state or every
possible tangent direction.  The same source studies $\mathcal{N}=4$ SYM and
finds vanishing curvature for the analyzed $1/4$-BPS sectors at generic
  coupling \cite{chen2026berry}.  In both cases the conclusion is a statement about the matrix
curvature of a degenerate bundle.  It is not the claim that the corresponding
Hamiltonian is integrable, nor that a vanishing curvature proves the absence
of all forms of quantum chaos.

The contrast is with the paper's black-hole-motivated examples.  In
$\mathcal{N}=2$ super-JT gravity, the Berry-curvature calculation gives
random-matrix-like curvature statistics; in the $\mathcal{N}=2$ SYK model,
explicit numerical calculations give the same qualitative behavior
\cite{chen2026berry}.  The word ``like'' matters.  The statement concerns the
statistics of the curvature matrix after the specified deformation and
symmetry prescription.  It does not say that the entries form an independent
Gaussian ensemble, that the distribution is Haar invariant, or that the
result is insensitive to the choice of moduli coordinates.  The gauge-invariant
content is carried by conjugation-invariant functions of $F$, while the
observable distribution may retain correlations between tangent directions.

This distinction is also visible in the relation between monotone and
fortuitous BPS states.  Supersymmetric SYK constructions and their protected
sectors are described in Refs.~\cite{fu2017susy,maldacenastanford2016}; the
finite-size appearance of additional relations is discussed in
Ref.~\cite{chang2024fortuity}.  The label ``fortuitous'' is a statement about
which BPS states exist at a given value of the number of degrees of freedom.
It is not itself a proof of random curvature.  Conversely, a monotone state
is not defined by $F=0$.  The curvature comparison in
Ref.~\cite{chen2026berry} supplies the additional calculation needed to
separate the two response behaviors.

\subsection{The \texorpdfstring{$\mathcal{N}=2$}{N=2} SYK and super-JT comparison}

The SYK model is useful here because it supplies a controlled large-$N$
setting with a well-developed low-energy description.  The random
spin-interaction precursor was introduced by Sachdev and Ye
\cite{sachdevye1993}.  The modern SYK analysis identifies an emergent
near-conformal regime and a reparametrization soft mode
\cite{maldacenastanford2016}; supersymmetric variants and their protected
sectors are constructed in Ref.~\cite{fu2017susy}.  These facts establish why
SYK is a useful model for testing a geometric diagnostic.  They do not turn an
SYK calculation into a theorem about higher-dimensional black holes.

The low-energy gravity comparison is similarly conditional.  Nearly-AdS$_2$
 dynamics is governed by a Schwarzian sector in the standard treatment
\cite{maldacenastanfordyang2016}, and the Schwarzian path integral and its
correlators have controlled formulations
\cite{stanfordwitten2017,mertens2017schwarzian}.  The direct statement relevant
to this paper is narrower: the $\mathcal{N}=2$ super-JT calculation in
Ref.~\cite{chen2026berry} produces random-matrix-like Berry curvature for the
protected states considered there.  Neither this calculation nor the SYK
numerics establishes that all supersymmetric black-hole microstates have a
universal curvature distribution.

The usual ETH language remains useful as a comparison of observables, not as a
replacement for the geometric calculation.  ETH concerns matrix elements of
operators in energy eigenstates, whereas the present object is a connection on
an exactly degenerate fiber.  SYK ETH studies
\cite{sonner2017syketh} and gravitational moment constructions
\cite{pollack2020gravityeth} help relate the two settings, but neither source
proves a Berry-curvature ETH statement.  The difference is especially
important at zero temperature, where a protected energy is fixed by a charge
constraint while the state bundle may still have nontrivial holonomy.

\subsection{A bounded relation to ``It from ETH''}

The phrase ``It from ETH'' refers here only to the construction of
Ref.~\cite{itfrometh2025}, which studies multi-interval entanglement and
replica-wormhole structures in a large-$c$ BCFT ensemble.  That work is a
relevant conceptual neighbor because it connects statistical information in a
boundary ensemble to geometric data in a gravitational description.  It does
not calculate the non-Abelian Berry curvature of a BPS state bundle, and it
does not validate the finite-channel law tested below.  We therefore use the
reference only to delimit a possible future bridge between state-bundle
geometry, entanglement data, and replica constructions.  No author attribution
other than the official record is made here, and no unresolved name is inserted
into the bibliography.

\subsection{Claim boundary for the present manuscript}

The direct evidence can be summarized without turning a model comparison into
a universal black-hole statement:
\begin{enumerate}
  \item the ordinary spectrum is silent inside an exactly degenerate BPS
  fiber;
  \item the cited D1/D5 and $\mathcal{N}=4$ SYM calculations provide
  non-random or vanishing curvature in the stated smooth/horizonless sectors;
  \item the cited $\mathcal{N}=2$ SYK and super-JT calculations provide
  random-matrix-like curvature in the stated black-hole-motivated models;
  \item the $\mathcal{N}=2$ SYK moduli-space calculation provides quantized
  Chern data, including exponentially large first Chern numbers in specified
  large-$N$ sectors \cite{chen2026berry}; and
  \item these results motivate tests of geometric response in additional
  protected systems, but do not establish that black holes satisfy an
  asymptotic or model-independent geometric-response law, ordinary ETH, or a
  thermalization law.
\end{enumerate}

This boundary is the one used in the remainder of the paper.  In particular,
the black-hole language identifies a falsifiable target class of models; it is
not used as an assumption that supplies the missing statistical closure.
